\documentclass[conference]{IEEEtran}

\usepackage[utf8]{inputenc}
\usepackage[T1]{fontenc}
\usepackage[spanish,english]{babel}
\usepackage{cite}
\usepackage{amsmath,amssymb,amsfonts}
\usepackage{graphicx}
\usepackage{float}
\usepackage{placeins}
\usepackage{booktabs}
\usepackage{textcomp}
\usepackage{url}
\usepackage{hyperref}



\begin{document}

\title{Sistema de detección de phishing basado en Machine Learning}

\author{John Jairo Riaño Martinez\
\thanks{Departamento de Ingeniería de Sistemas e Industrial, Universidad Nacional de Colombia, Bogotá, Colombia. Correo: jorianom@unal.edu.co. Asignatura: Ciberseguridad.}}

\maketitle

\begin{abstract}
This technical report describes the development of a phishing URL detection system based on machine learning. The project covers the full pipeline, from data collection and cleaning to feature engineering, model training, evaluation, and deployment. The selected model is exposed through a FastAPI backend and consumed from a simple web interface. The document focuses on the practical design of the system, the interpretability of URL-based features, and the experimentally verified results obtained from the notebook. The final solution is designed for fast inference and straightforward operational use.
\end{abstract}

\begin{IEEEkeywords}
phishing detection, machine learning, feature engineering, FastAPI, cybersecurity, URL analysis
\end{IEEEkeywords}

\section{Introducción}
El phishing sigue siendo uno de los vectores de ataque más efectivos en ciberseguridad porque combina ingeniería social con infraestructura web diseñada para suplantar servicios legítimos. En este contexto, la URL no es solo un identificador técnico del recurso, sino también una fuente de señales que permiten inferir comportamiento malicioso antes de que el usuario interactúe con la página.

Los enfoques basados en listas negras siguen siendo útiles, pero reaccionan después de que un dominio ya fue identificado. Cuando los atacantes registran dominios nuevos de forma continua, la velocidad de respuesta de estos mecanismos se vuelve insuficiente. Por esta razón, el proyecto adopta un enfoque de Machine Learning que clasifica URLs a partir de características léxicas y estructurales, sin depender de navegación automática ni de consultas externas en tiempo real.

El objetivo del sistema es construir una solución completa y utilizable: un notebook de experimentación, un modelo persistido, una API de inferencia y una interfaz web ligera para consulta rápida.

\section{Metodología y desarrollo}
\subsection{Datos y recolección}
La base de datos se construyó a partir de archivos ubicados en \texttt{data/raw/}, con fuentes orientadas a phishing y dominios legítimos. En el repositorio aparecen conjuntos derivados de PhishTank, URLhaus, OpenPhish, ISCX, Umbrella y Tranco. Esta combinación permite entrenar un clasificador con ejemplos de comportamiento malicioso y benigno.

Las fuentes principales se usan con dos propósitos: por un lado, enriquecer la clase phishing con URLs reportadas por sistemas de inteligencia de amenazas; por otro, incorporar URLs legítimas representativas para reducir sesgos hacia dominios triviales o poco variados.

\subsection{Preprocesamiento}
El preprocesamiento normaliza la cadena de entrada, estandariza el esquema de la URL y extrae componentes como subdominio, dominio y sufijo con \texttt{tldextract}. También se limpian elementos redundantes para reducir ruido en las variables derivadas y evitar inconsistencias entre fuentes.

\subsection{Ingeniería de características}
El sistema transforma cada URL en un vector numérico construido localmente. Las características observadas en el notebook y en \texttt{api/index.py} se agrupan en cuatro bloques.

\begin{itemize}
    \item \textbf{Conteos y longitudes:} longitud total de la URL, longitud del hostname, path y query; número de puntos, guiones, dígitos, parámetros y otros separadores.
    \item \textbf{Relaciones y proporciones:} razón de dígitos, proporción de caracteres especiales y secuencias numéricas largas.
    \item \textbf{Indicadores binarios:} presencia de IP en la URL, símbolo \texttt{@}, acortadores, TLD sospechoso, punycode, doble barra en el path y palabras clave asociadas con inicio de sesión o actualización de cuenta.
    \item \textbf{Estructura del dominio:} longitud del dominio, longitud del sufijo, número de subdominios, entropía del hostname y patrones de consonantes.
\end{itemize}

Este diseño prioriza interpretabilidad: cada variable responde a una intuición clara sobre patrones observables en URLs usadas en campañas de phishing.

\subsection{Modelado}
Durante el entrenamiento se compararon varios modelos supervisados, incluyendo \texttt{RandomForestClassifier}, \texttt{ExtraTreesClassifier} y \texttt{GradientBoostingClassifier}. La selección final se basó en el desempeño en validación y en la estabilidad frente a datos no vistos. El pipeline ganador se serializó con \texttt{joblib} para garantizar que el mismo preprocesamiento y la misma lógica de predicción se reutilicen en inferencia.

La versión persistida también incluye el umbral operativo que maximiza F1 en el conjunto de validación, lo que permite ajustar la decisión binaria sin modificar el modelo base.

\subsection{API del modelo}
La API se implementó en \texttt{api/index.py} con \textbf{FastAPI}. El flujo de inferencia es directo: el cliente envía una URL en formato JSON, la API normaliza la entrada, extrae las mismas características del notebook y las envía al pipeline serializado. La respuesta retorna la probabilidad de phishing, el veredicto final y señales explicativas basadas en reglas heurísticas.

Importante: tanto en el notebook (variable \texttt{TRANCO\_TOP10K} usada por la clase \texttt{PhishingDetector}) como en el backend (lista de dominios cargada por la API) se incorpora una regla ligera basada en la "Tranco Top\,-10k". La lista Tranco es una referencia pública de los sitios más visitados que aquí se usa como proxy de "dominios ampliamente conocidos". Si el dominio raíz de la URL figura en esa lista, el detector puede anular la clasificación a \texttt{LEGÍTIMA} (override) para reducir falsos positivos sobre sitios de alto tráfico.

Esta heurística es una medida práctica de filtrado rápido —no una verificación de seguridad exhaustiva—; su propósito es minimizar alarmas en dominios muy populares, y por tanto debe entenderse como un atenuador de la decisión del clasificador, no como un sustituto de validaciones adicionales.

\subsection{Interfaz web}
La interfaz web se encuentra en \texttt{docs/index.html}. Su función es ofrecer una capa de consumo simple donde el usuario ingresa una URL y visualiza el resultado sin complejidad adicional. La interfaz presenta el veredicto, la probabilidad y las señales detectadas en tarjetas de estado, con un diseño pensado para lectura rápida y operación práctica.

La interfaz está desplegada en \url{https://jorianom.me/project_cyber/} y puede probarse directamente en línea.

\begin{figure}[H]
\centering
\includegraphics[width=0.95\linewidth]{figures/frontend_web.png}
\caption{Interfaz web — versión desplegada.}
\label{fig:frontend_web}
\end{figure}

\section{Resultados}
La evaluación del modelo ganador en el notebook produjo una accuracy de 0.9088 sobre 19{,}977 ejemplos de prueba. El classification report mostró una media macro de precision, recall y F1 de 0.9091, 0.9088 y 0.9088, respectivamente. Además, el umbral operativo óptimo encontrado fue 0.47, con F1 = 0.9088, recall = 0.9060 y precision = 0.9117. El notebook también reporta ROC-AUC = 0.9732 y average precision = 0.9759.

La matriz de confusión registrada en la salida del notebook fue la siguiente:

\begin{table}[H]
\caption{Matriz de confusión del modelo ganador}
\label{tab:confusion}
\centering
\begin{tabular}{lcc}
\toprule
 & \textbf{Pred. Legítima} & \textbf{Pred. Phishing} \\
\midrule
\textbf{Real Legítima} & 9207 & 770 \\
\textbf{Real Phishing} & 1052 & 8948 \\
\bottomrule
\end{tabular}
\end{table}
\FloatBarrier

La Fig.~\ref{fig:evaluation} resume la evaluación global del modelo, mientras que la Fig.~\ref{fig:importance} muestra las variables con mayor peso en el clasificador. La Fig.~\ref{fig:threshold} ilustra el ajuste del umbral de decisión.

En términos de interpretación, el comportamiento del clasificador confirma que las variables estructurales de la URL aportan una señal útil para distinguir entre enlaces legítimos y enlaces usados en campañas de phishing.

\begin{figure}[H]
\centering
\includegraphics[width=0.95\linewidth]{figures/evaluation.png}
\caption{Evaluación del modelo: matriz de confusión, curva ROC y curva Precision-Recall.}
\label{fig:evaluation}
\end{figure}

\begin{figure}[H]
\centering
\includegraphics[width=0.82\linewidth]{figures/feature_importance.png}
\caption{Importancia de variables del modelo ganador.}
\label{fig:importance}
\end{figure}

\begin{figure}[H]
\centering
\includegraphics[width=0.82\linewidth]{figures/threshold.png}
\caption{Comportamiento de precision, recall y F1 en función del umbral de decisión.}
\label{fig:threshold}
\end{figure}

\section{Discusión}
El principal valor del sistema es su simplicidad operacional. La inferencia no depende de análisis de contenido ni de consultas remotas, por lo que el tiempo de respuesta es bajo y la arquitectura es apropiada para consumo en línea. Además, el uso de variables explícitas facilita explicar por qué una URL fue clasificada como sospechosa.

La comparación entre modelos sugiere que los ensambles de árboles se adaptan bien al problema, probablemente porque capturan relaciones no lineales entre señales léxicas y estructurales. Sin embargo, la robustez del sistema sigue limitada por el tipo de datos disponibles: si una URL legítima presenta patrones poco convencionales, el modelo puede elevar su probabilidad de phishing; de forma equivalente, un ataque muy bien disfrazado puede evadir las reglas heurísticas.

\section{Limitaciones}
El enfoque presenta limitaciones importantes. Primero, el sistema depende de un dataset finito y requiere reentrenamiento periódico para seguir el ritmo de nuevas campañas. Segundo, las features son principalmente léxicas; por tanto, el sistema no inspecciona el contenido real de la página ni su comportamiento dinámico. Tercero, los falsos positivos y falsos negativos siguen siendo posibles cuando la URL se aleja de los patrones aprendidos.

\section{Conclusiones}
Se construyó un sistema completo para detección de phishing que integra preparación de datos, ingeniería de características, entrenamiento, evaluación, persistencia e inferencia. El aporte más importante no es la complejidad del modelo, sino la coherencia entre el notebook, el backend y la interfaz web. Esa consistencia hace que el sistema sea reproducible y fácil de operar.

Los resultados del notebook muestran un desempeño sólido para un enfoque basado en URL: accuracy de 0.9088, ROC-AUC de 0.9732 y average precision de 0.9759. En conjunto, estos valores respaldan la utilidad del sistema como base práctica para detección rápida de phishing y como punto de partida para mejoras futuras.

\begin{thebibliography}{00}
\bibitem{colab}
J.~Riaño, ``Project Cyber — Jupyter notebook (Colab),'' 2026. [Online]. Available: \url{https://colab.research.google.com/github/jorianom/project_cyber/blob/main/ml_pishing.ipynb}

\bibitem{github}
Jorianom, ``Project Cyber repository,'' 2026. [Online]. Available: \url{https://github.com/jorianom/project_cyber}

\bibitem{course}
Cybersec Course, ``Material de referencia del curso,'' 2026. [Online]. Available: \url{https://sites.google.com/view/cybersec-2026-1}
\end{thebibliography}

\end{document}